\documentclass[11pt]{article}
\usepackage[margin=1in]{geometry}
\usepackage{amsmath,amssymb,amsthm,booktabs}
\usepackage{needspace}
\usepackage[hidelinks]{hyperref}
\setlength{\emergencystretch}{2em}
\newtheorem{theorem}{Theorem}
\newtheorem{proposition}[theorem]{Proposition}
\newtheorem{corollary}[theorem]{Corollary}
\newcommand{\F}{\mathbb F}
\newcommand{\Z}{\mathbb Z}
\newcommand{\mr}{\operatorname{mr}}
\newcommand{\rank}{\operatorname{rank}}
\newcommand{\im}{\operatorname{im}}
\title{Rank-two affine certificates are exactly three-sheeted cover homology}
\author{Internal research working proof}
\date{17 November 2026}
\begin{document}
\maketitle

The rank-two-form part of the original quadratic-colouring method
has an exact cohomological description, including arbitrarily large
radicals. Fix an onto transposition colouring
$\rho:\Gamma=\pi_1(S^3\setminus K)\to S_3=GL_2(\F_2)$, and let
$Y_\rho$ be its simple three-sheeted branched cover. Write
$h=b_1(Y_\rho;\F_2)$. For $k\ge2$, a transitive affine
representation of the specified kind, with form rank two and
dimension $k$, projecting to $\rho$, exists exactly when
\[
                         h\ge k-1.
\]
Thus when $h>0$ its best bound is $\mr(K)\ge h+2$, and its largest
dimension is $h+1$. When $h=0$ no such representation exists.
This supplies both a construction and an exact stopping condition
for this part of the user's conditional theorem.

There is also a quotient-count version: the number of distinct
$S_4$ transposition quotient kernels above this fixed $S_3$ kernel
is $2^h-1$. Two such $S_4$ kernels on the same knot group therefore
force $\mr(K)\ge4$ and construct a transitive eight-point
rank-two-form certificate. They must share the same $S_3$ quotient;
kernels on different knots or different lower quotients cannot be
combined in this assertion.

The accompanying complete 293-knot test has no new sharp case.
The uniform arguments below are working proofs, unformalized and
externally unreviewed, without a novelty or full MRC claim.

\section{A general radical-lifting lemma for knot groups}

For a representation on $M$, let $Z^1_{\rm ad}(\Gamma,M)$ consist
of cocycles $c$ with
$c(m)\in\im(\rho(m)-I)$ for every meridian $m$. In a meridional
presentation whose meridian differences have rank one, evaluation
gives one scalar per generator and injects this space into
$k^d$ for any $d$ generating meridians. Consequently
$\mr(K)\ge\dim Z^1_{\rm ad}(\Gamma,M)$.

We use two elementary knot-group facts, with trivial coefficients:
\[
 H^1(\Gamma,R)\xrightarrow{\ c\mapsto c(m)\ }R\text{ is an
 isomorphism},\qquad H^2(\Gamma,R)=0.
\]
Indeed, the exterior has $H_1=\Z$ generated by a meridian and
$H_2=0$. Attach cells of dimension at least three to make a
classifying space. This cannot create second homology, so
$H_2(\Gamma;\Z)=0$. The universal coefficient theorem gives the
claims for every vector space $R$ with trivial action. No
asphericity assumption is needed for this argument.

\Needspace{9\baselineskip}
\begin{proposition}\label{prop:radical}
Suppose $0\to R\to V\xrightarrow{\pi}W\to0$ is a sequence of
$k\Gamma$-modules with $R$ trivial. For every meridian, suppose
$\rho_V(m)-I$ and $\rho_W(m)-I$ both have rank one and their image
lines map isomorphically under $\pi$. Then
\[
 \pi_*:Z^1_{\rm ad}(\Gamma,V)\longrightarrow
             Z^1_{\rm ad}(\Gamma,W)
\]
is an isomorphism over every field $k$.
\end{proposition}
\begin{proof}
The long exact cohomology sequence and $H^2(\Gamma,R)=0$ lift
the class of any admissible $W$ cocycle to $V$. Lifting a vector
and adding its principal cocycle makes the projection equal to
the specified cocycle, not merely its class. At one meridian $m$,
subtract the unique homomorphism $\Gamma\to R$ which puts its
value in the required image line. Such a homomorphism exists by
meridian evaluation on $H^1(\Gamma,R)$.

Homogenize this cocycle to the representation
$g\mapsto\left(\begin{smallmatrix}\rho_V(g)&c(g)\\0&1\end{smallmatrix}\right)$.
Its difference from identity at $m$ has rank one. All meridians
are conjugate (or inverses of conjugates), so this remains true
at every meridian, proving admissibility. For injectivity, a
cocycle projecting to zero is an $R$-valued homomorphism. At $m$
its value lies in $R\cap\im(\rho_V(m)-I)=0$, hence it is zero.
\end{proof}

In particular, for an $S_4$ transposition colouring over $\F_2$,
the deleted permutation module $D\subset\F_2^4$ has a fixed line
$R=\langle(1,1,1,1)\rangle$, and $D/R=W$ is the natural
two-dimensional module for $S_4/V_4=S_3$. The image lines survive
in the quotient, so its binary scalar root-lift dimension is
exactly the dimension for the induced $S_3$ colouring. This
compares particular coefficient modules, not their integral
Smith forms or arbitrary other characteristics.

\section{Normal form and necessity}

On $W=\F_2^2$, put
\[
 B(w,v)=w_1v_2+w_2v_1,\qquad
 Q_0(w)=w_1+w_2+w_1w_2.
\]
The quadratic form $Q_0$ equals one at every nonzero vector.
The three transpositions of $S_3$ act as
$T_v(w)=w+B(w,v)v$, $v\ne0$.
For the fixed colouring let $L=Z^1_{\rm ad}(\Gamma,W)$, expressed
as scalar arrays $a_i$ with $c(m_i)=a_i v_i$. The principal arrays
are $a_i=B(s,v_i)$, $s\in W$, and have dimension two. Set
\[
                    \ell=\dim L,\qquad H=L/B^1,\qquad h=\ell-2.
\]
This $H$ embeds in ordinary $H^1(\Gamma,W)$; it need not be all
ordinary cohomology. Equivalently it is the first cohomology
after imposing meridian squares, in this coefficient module.

Let $V=W\oplus R$, $R=\F_2^q$, with form $B$ pulled back from
$W$. Every affine meridian of the specified original kind has
the normal form
\begin{equation}\label{eq:normal}
 (w,r)\longmapsto
 \bigl(w+(B(w,v_i)+e_i)v_i,\quad
       r+(B(w,v_i)+e_i)z_i\bigr),
 \qquad z_i\in R, e_i\in\F_2.
\end{equation}
The form radical is fixed pointwise by the linear parts.
Transitivity forces the quotient roots to span $W$: otherwise
the quotient by their span is fixed. Pure translations cannot
give an exception: if one linear meridian part is identity they
all are, and a cyclic translation image cannot be transitive
in dimension at least two. Thus the quotient linear image is
$S_3$.

The array $e$ belongs to $L$ by the affine action on $W$.
Each coordinate array $z^j=(z_i^j)_i$ also belongs to $L$:
the corresponding linear extension by a trivial line is dual
to the natural cocycle through $B$. This follows as well by
inserting its lower-left blocks $z_i^j B(-,v_i)$ in the relators.

\Needspace{6\baselineskip}
\begin{proposition}\label{prop:independence}
Transitivity in \eqref{eq:normal} forces
$[e],[z^1],\ldots,[z^q]$ to be independent in $H$.
\end{proposition}
\begin{proof}
A nonzero dependence would give $\alpha\in R^*$,
$\beta\in\F_2$, not both zero, and $s\in W$ such that
\[
          \alpha(z_i)+\beta e_i+B(s,v_i)=0\quad\text{for all }i.
\]
The function
\[
             F(w,r)=\alpha(r)+\beta Q_0(w)+B(s,w)
\]
is then invariant under every meridian. To check this, write
$a=B(w,v_i)$. The change of $Q_0$ is $e_i(a+1)$, so the total
change is
\[
 [\alpha(z_i)+\beta e_i+B(s,v_i)]a
       +e_i[\alpha(z_i)+\beta+B(s,v_i)]=0.
\]
The constant term vanishes because $e_i^2=e_i$.
If $\alpha\ne0$ the function is nonconstant in $r$; otherwise
$\beta=1$ and its nonzero polar form makes it nonconstant in
$w$. This contradicts transitivity.
\end{proof}

Since $k=q+2$, this proves $h\ge q+1=k-1$. It treats every
radical dimension, without enumerating affine groups.

\section{Sufficiency, explicit construction and image order}

Conversely choose independent classes $[e],[z^1],\ldots,[z^q]$
in $H$, represented by arrays in $L$. Their linear block
extensions define an action on $V=W\oplus R$ with meridians
\[
 A_i(w,r)=\bigl(T_{v_i}w,\ r+B(w,v_i)z_i\bigr).
\]
Its meridian image line is spanned by $(v_i,z_i)$ and maps
nontrivially to $W$. Proposition~\ref{prop:radical} lifts the
admissible cocycle $e$ uniquely to $V$; hence its meridian values
are $e_i(v_i,z_i)$. This proves that the explicit formulas
\eqref{eq:normal} satisfy all knot-group relations. An arbitrary
finitely presented group cannot replace the knot group in this
step: the vanishing and meridian-evaluation hypotheses matter.

It remains to prove transitivity, not just lack of a fixed point.
Let $N=\ker\rho$. The natural $S_3$ module $W$ is simple with
$\operatorname{End}_{S_3}(W)=\F_2$, and $H^1(S_3,W)=0$.
For the last assertion, average over the normal odd-order $C_3$;
its invariants on $W$ vanish. Inflation--restriction therefore
injects $H^1(\Gamma,W)$ into
$\operatorname{Hom}(N,W)^{S_3}$.

The chosen classes restrict to independent equivariant maps.
Their joint image is the full module $W^{q+1}$: a proper
submodule of this semisimple module is annihilated by a nonzero
$S_3$-map to $W$, which would give a dependence. This uses
semisimplicity of this direct sum of simples, not semisimplicity
of the whole modular group algebra.

On $N$, write the affine transformations as triples $(Z,t,s)$,
acting by $(w,r)\mapsto(w+t,r+Zw+s)$, with
$Z\in\operatorname{Hom}(W,R)$, $t\in W$, $s\in R$. Multiplication
and commutation give
\[
 (Z,t,s)(Z',t',s')=(Z+Z',t+t',s+s'+Zt'),
 \qquad [g,g']=(0,0,Zt'+Z't).
\]
The full joint image in $(Z,t)$ supplies a $Z$ supported in any
chosen radical coordinate and a $t'$ with $Zt'\ne0$.
These commutators supply every radical translation. Thus the
image on $N$ contains all triples and acts transitively on $V$.

\begin{theorem}\label{thm:equivalence}
For the fixed onto $S_3$ colouring and $k\ge2$, the original
rank-two-form transitive affine certificate in dimension $k$
exists if and only if $h\ge k-1$. Every such transitive image
is the full affine group with linear parts fixing $R$ pointwise
and inducing $S_3$ on $V/R$; its order is
\[
                     24\,8^{k-2}.
\]
For maximal dimension $k=h+1$, the original bound is $h+2=\ell$.
\end{theorem}

The order follows from the six choices on $W$, $2^{2q}$ choices
of $Z$ and $2^{q+2}$ translations. Necessity already supplies
the independence used above, so the image statement holds for
every transitive representation of the specified kind, not only
the constructed maximal ones.

\section{Cover homology and quotient-kernel consequences}

The retained branched-cover theorem identifies the scalar lift
space of a degree-$r$ transposition colouring over any field:
$\dim Z^1_{\rm ad}=r-1+b_1(Y;k)$. Its proof eliminates the fixed
sheet and squared-meridian cycles in the lifted presentation.
For $r=3$ it gives $h=b_1(Y_\rho;\F_2)$ and the description
announced at the start. There is no claim that this bounds
ordinary double-cover group rank or detects every knot.

For the quotient-count version, each $c\in L$ gives an affine
action on $W$, hence a lift of $\rho$ to
$AGL_2(\F_2)=S_4$. Its meridians are transpositions: an
admissible affine transvection has two fixed points on four
points. Principal cocycles give conjugate complements $S_3$.
Every nonprincipal class gives the full $S_4$: the translation
kernel is either zero or all of the irreducible $W$; in the
zero case $H^1(S_3,W)=0$ makes the cocycle principal.

Conversely every $S_4$ transposition quotient above $\rho$ is
such an affine lift. Automorphisms of $S_4$ are inner; those
inducing the identity on $S_3=S_4/V_4$ act here by the four
translations. Thus quotient kernels above the fixed $S_3$
kernel correspond exactly to the nonzero classes in $H$.
There are $2^h-1$. In particular,
\[
 \mr(K)\le d\quad\Longrightarrow\quad
 \#\{S_4\text{ kernels above this }S_3\text{ kernel}\}
                         \le 2^{d-2}-1.
\]
Two distinct such kernels give $h\ge2$, a transitive
eight-point certificate and lower bound four. The third
nonzero class is then forced, rather than counted across a
different knot or a different lower quotient. This is a
relative version of the user's joint-quotient idea, with no
priority claim.

Proposition~\ref{prop:radical} gives a further mod-two identity.
For an $S_4$ transposition colouring, let $Y_4$ be its simple
four-sheeted cover and $Y_3$ the cover of its $S_3$ quotient.
Equality of the two admissible dimensions gives
\[
                   b_1(Y_3;\F_2)=1+b_1(Y_4;\F_2).
\]
No integral homology or all-characteristic analogue is asserted.

\section{The complete 293-knot panel}

All names, PDs and four-seed certificates are from the existing
pinned open-list panel, in the same order. The new frozen test
exhausts connected symmetric transposition colourings up to
sheet relabelling. There are respectively $1,13,39,25$ seed
orbits in degrees $2,3,4,5$. All accepted representatives have
exact integer scalar matrices and replayable Smith reductions.
The totals are
\[
\begin{array}{c|rrrr}
\text{degree}&2&3&4&5\\\hline
\text{quotient representatives}&293&26&21&0.
\end{array}
\]
For all 26 degree-three bases, the binary lift dimension is two
on five and three on 21. Thus $h=0$ or one; no target reaches
the sharp-four threshold. Exactly 21 knots admit the original
criterion in dimension two, giving bound three. The other 272
admit no representation satisfying the original nonzero-form
criterion in any dimension. Indeed its own bound and the
upper certificate force $k\le3$, hence form rank two; the
equivalence above then applies. None of the 293 admits a sharp
original quadratic certificate. This is an exclusion of this
method, not a proof that their meridional ranks are three.

Across every symmetric degree and every characteristic, the
maximum scalar bounds are one on six knots, two on 251 and
three on 36. Degrees at least six are excluded by the connected
graph bound: four generating transpositions connect at most
five sheets. Independently, all 170 retained $A_5$ cover
quotient kernels have inner deck action. An outer deck action
would extend to an $S_5$ transposition quotient and force sharp
four; the absence agrees exactly with the direct degree-five
enumeration. The known $T(4,5)$ permutation control has the
outer extension and verifies that this gate can pass.

\section{Completing the binary affine three-dimensional base gap}

Every rank-one invertible meridian matrix over $\F_2$ in
dimension three is a transvection. Its image group is normally
generated by that involution. Independent enumeration of all
179 subgroups of $GL_3(\F_2)$, with closure of the supplied
lattice under adjoining every element, leaves only image
types $C_2,S_3,S_4,GL_3(\F_2)$. Their numbers of embedded
subgroups are $21,28,14,1$. The corresponding modules are
handled as follows.

For $C_2$ all meridians are the same involution and the
admissible space has dimension one. For $S_3$, the module is
$W\oplus\F_2$ and all root lines lie in $W$, giving the same
space $L$. The two types of $S_4$ subgroup are dual. One has
the deleted module with a fixed line, covered by
Proposition~\ref{prop:radical}. In the other, every root lies
in a two-dimensional $S_3$ submodule, so admissible cocycles
already take values there. Their spaces are again $L$.

For the full $GL_3$ image, restriction to the double-cover
group is onto since the image has no subgroup of index two.
The deck action must be inner. The retained complete GL3
audit already checks both natural and dual actions for all
29 inner-deck kernels, with root dimension three in all 58
cases. Those results are reused by hash, without repeating
the completed quotient search. The subgroup classification
and the new $S_3/S_4$ tests therefore complete the proper-image
gap: every three-dimensional binary linear base on every
one of the 293 targets has admissible dimension at most three.
No homogenized affine extension of such a base can give the
missing sharp-four certificate. Other fields, dimensions,
nonlinear mechanisms and representation components remain open.

\section{A nonlinear three-sheeted cover follow-up}

The retained general cover bound is
$\mr(K)\ge2+\rank\pi_1(Y_3)$. Thus one nonabelian finite
quotient of a degree-three cover group would give sharp four,
even when its homology is cyclic. All 26 covers were frozen
for this different test. A graph presentation kills fixed-sheet
meridian loops, identifies the exchanged-sheet edges as inverses,
and gauges two spanning-tree edges to identity. Every knot
relator is then lifted from all three starting sheets.

Explicit one-occurrence Tietze eliminations prove that 19 of
these cover groups are cyclic: seventeen of order two, one
of order four and one trivial. Seven presentations remain on
two generators. Exhaustive evaluation of every pair in
$S_3,A_4,D_{10},D_{14},A_5$ finds no nonabelian image for them;
here $D_{10},D_{14}$ have orders ten and fourteen. This does
not prove those seven groups cyclic.

A further frozen Todd--Coxeter proposal with at most 5,000
cosets per presentation returns cyclic permutation images for
three and reaches the cap on four. The returned actions are
independently checked only as quotients; their cyclicity is
not adopted as a proof of source cyclicity. The remaining
cover-group ranks and possible nonabelianity are unresolved.
No three-bridge knot presentation follows from a cyclic cover
group alone. The two exact replay receipts retain every
presentation, elimination, finite test and cap outcome.

\section{Verification and limits}

\texttt{rank\_two\_equivalence\_verification.json} seals the
plans, complete panel and controls. The replay uses full tuple
orbit marking, fixed-trace permutation closure, reverse affine
point evaluation and independent Smith-step replay. A separate
$PSL_2(\F_5)$ model checks the deck classification. Source PD
parsing and the older Smith replay routine are shared.

Three actual known pretzel diagrams $P(2,3,3)$,
$P(2,3,3,3)$ and $P(2,3,3,3,3)$ are constructed with pinned
Spherogram 2.4.1. Their maximal binary scalar dimensions are
three, four and five. Every cocycle-basis pair satisfies the
quadratic affine relators; independent classes produce
transitive actions on four, eight and sixteen points, of
orders 24, 192 and 1536. Dependent choices preserve the
displayed nonconstant quadratic invariant. The induced $S_4$
modules have exactly the same binary lift dimensions. These
are known sharp pretzel controls covered by the retained
published Coxeter/Montesinos results, not new open cases.

The computations verify finite hypotheses and examples, not
the uniform proofs or full MRC. Actual ranks and bridge
equalities for the 293 open-list targets are unchanged.
No formal verification, specialist review or general finite
quotient detection theorem is claimed. Full classical MRC and
the broader source/priority audit remain open.
\end{document}
